% !TITLE 早发白帝城
% !AUTHOR 李白
% !DYNASTY 唐
% !ORDER 16

\begin{poem}
朝辞白帝彩云间\textsuperscript{1}，千里江陵\textsuperscript{2}一日还。
两岸猿声啼不住\textsuperscript{3}，轻舟已过万重山。
\end{poem}

\begin{annotations}
\notetext{1}{白帝：白帝城，在今重庆奉节，长江三峡的起点。城在高山上，云雾缭绕，故称“彩云间”。}
\notetext{2}{江陵：今湖北江陵，距白帝城约一千二百里。顺江而下水流极快，所以“一日还”虽是夸张，但也有现实依据。}
\notetext{3}{啼不住：叫个不停。郦道元《水经注》说“巴东三峡巫峡长，猿鸣三声泪沾裳”，猿声本为悲音，但李白此时遇赦东归，猿声成了轻舟疾驰的背景音。}
\end{annotations}

\begin{appreciation}
这是李白晚年写得最快、最轻的一首诗，也是中国文学史上写“归心似箭”写得最好的绝句之一。

先说背景，这二十八个字背后站着一个死里逃生的人。李白晚年卷入一场叛乱，被流放夜郎——那时的夜郎在今天贵州一带，对一个快六十岁、一身病痛的老人来说，基本等于等死。他沿长江一路往西南走，走到白帝城，忽然传来大赦的消息。上一刻还在赴死，下一刻无罪放归。你设身处地想一想，就明白他这一天的船划得有多快。

朝辞白帝彩云间。白帝城在瞿塘峡口的高山上，云绕雾罩，像在彩云里。人高兴的时候，看什么都像仙境。千里江陵一日还。白帝到江陵一千二百里，顺流而下，一天就到。郦道元《水经注》写过这段水路：“朝发白帝，暮到江陵”，说是“虽乘奔御风，不以疾也”——骑着快马、驾着疾风，也没它快，不是瞎吹，有现实根据。但“还”字才是诗眼：是回去。是被赦免的人回自己该回的地方，每一桨下去都是回家。

两岸猿声啼不住，轻舟已过万重山。三峡的猿声本是出名的悲音，《水经注》说“巴东三峡巫峡长，猿鸣三声泪沾裳”。同一条江，同一片猿声，从前的人听出眼泪，李白只听见风从耳边过。不是猿声变了，是心情变了——一个人从地狱门口折返，世上再没有什么声音能把他留住。轻舟已过万重山，那万重山是他这一辈子的坎儿，都过去了。

我们读过他年轻时写的《渡荆门送别》，少年李白出蜀，顺着这条江第一次去看世界，“山随平野尽，江入大荒流”，满眼都是前程。几十年后，还是这条长江，他往回走，走成一首二十八个字的绝句。一条江，一个人，一出一回，就是一生。

这首诗哥哥希望你背下来。你以后会碰到一些觉得熬不过去的坎——考试、升学、为难的事，当时觉得天大的，回头看，不过万重山里的一重。都会过去的。真到了那一天，你也坐上你的那艘“轻舟”，就知道李白没骗人。
\end{appreciation}
